% !TeX root = ../tfg.tex
% !TeX encoding = utf8
%
%*******************************************************
% Capítulo 15. Conclusiones y trabajo futuro
%*******************************************************

\chapter{Conclusiones y trabajo futuro}\label{cap:conclusiones}

\section{Conclusiones}\label{sec:conclusiones}

% Revisar objetivo por objetivo (\cref{sec:objetivos}) si se ha alcanzado, y responder a las
% preguntas de investigación (\cref{sec:preguntas}).
\pendiente{conclusiones, objetivo por objetivo}

\section{Limitaciones}\label{sec:limitaciones}

La detección OOD protege frente a imágenes raras o desconocidas, pero no frente a cambios en el criterio diagnóstico. Para eso hace falta revisar y reetiquetar los datos periódicamente.

% Otras limitaciones a desarrollar:
% - Fuga por lámina: la evaluación mide generalización dentro de lámina (\cref{cap:fuga}).
% - Partición de entrenamiento y validación no reproducible (sin semilla).
% - La validación se usó para elegir la época: intercambiabilidad no asegurada (\cref{cap:conforme}).
% - Hipótesis del label smoothing no comprobada (habría que entrenar sin él).
% - Interpretación de STOOD-X pendiente de confirmar con los autores.
% - Límites teóricos de la detección OOD: aprendibilidad (Fang et al., 2022).
\pendiente{límites teóricos de la detección OOD: aprendibilidad (Fang et al.~\cite{fang2022})}
\pendiente{resto de limitaciones}

\section{Trabajo futuro}\label{sec:trabajo-futuro}

% - Explicabilidad y eliminación de sesgos con CRP (\cref{cap:xai-sesgos}).
% - Reentrenar con partición por paciente o validar con un conjunto externo de próstata.
% - Entrenar sin label smoothing para contrastar la hipótesis sobre MSP, MaxLogit y Energy.
\pendiente{trabajo futuro}

\endinput
